% Mirror of abstract in main.tex — do not edit here.

Multi-judge LLM evaluation panels are assumed to benefit from diversity, yet diversity's relationship to adversarial vulnerability is poorly understood. We decompose panel attack success into individual judge weakness ($\eta_{\max}$) and structural diversity ($K_{\mathrm{eff}}$) across 15 models, six attack types, two benchmarks, and panel sizes $K \in \{3, 5, 7\}$ (up to 6,435 panels). Consistent with adversarial ensemble theory but newly characterized in the LLM judge setting, $K_{\mathrm{eff}}$ is uniformly positively associated with vulnerability across all 12 conditions, though the diversity--capability confound ($r{=}0.70$--$0.81$) limits causal interpretation. The predictors' relative importance shifts with panel size: at $K{=}3$, both contribute; at $K{\geq}5$, $K_{\mathrm{eff}}$ becomes the dominant predictor while $\eta_{\max}$ sharply attenuates. Cinelli--Hazlett sensitivity analysis reveals that $K_{\mathrm{eff}}$'s robustness value does not exceed the known confound in 69\% of conditions (25/36), but fractional logit with clustered standard errors confirms the positive association in most conditions ($K{=}3$: 9/12; $K{=}5$: 11/12; $K{=}7$: 12/12). Aggregation rules influence predictability ($R^2$: threshold 0.83, Bayesian 0.09). The defense strategy should be informed by panel size: small panels benefit from screening for individual robustness; larger panels benefit from monitoring judge independence.
